\documentclass[]{article}

\usepackage{amsmath}
\usepackage{multirow}
\usepackage{natbib}
\usepackage{graphicx} 


\begin{document}

The performance of the Multi-Frequency Wiener Filter (MWF) and the standard Internal Linear Combination (ILC) was evaluated by applying them to simulated multi-frequency sky maps. We first assess the numerical challenges and the statistical properties of the input signals, then compare the reconstruction fidelity in the harmonic domain, and finally evaluate the scientific utility of the resulting Compton-$y$ maps through a galaxy cluster detection pipeline.

\subsection{Input signal properties and numerical considerations}

A primary challenge in this analysis is the vast dynamic range of the signal components across the six frequency channels. As illustrated in the left panel of Figure \ref{fig:image0}, the variance in the high-frequency channels (545 and 857 GHz), which are dominated by the Cosmic Infrared Background (CIB), is approximately five orders of magnitude greater than in the lower-frequency channels dominated by the CMB. The target tSZ Compton-$y$ signal itself has a variance many orders of magnitude smaller than any of the observed channels. This disparity necessitates the use of high-precision arithmetic and numerical stabilization techniques, such as the standardization of the target $y$-map (Figure \ref{fig:image0}, center and right panels), to ensure the covariance matrix used by the MWF is well-conditioned and invertible.

\begin{figure}[h!]
    \centering
    \includegraphics[width=0.8\textwidth]{../input_files/plots/step_1_diagnostic_plot_1_20260416_234622.png}
    \caption{The extreme dynamic range of the input signals, which presents a numerical challenge for the Wiener filter. The per-channel signal variance (left) shows that the high-frequency channels (545 and 857 GHz), dominated by the Cosmic Infrared Background (CIB), have a variance approximately five orders of magnitude larger than the lower-frequency channels. The distribution of the target tSZ Compton-$y$ signal is shown in its raw, small-valued form (center) and after standardization (right), a crucial pre-processing step that ensures the numerical stability of the covariance matrix inversion required by the filter.}
    \label{fig:image0}
\end{figure}

With the numerical challenges addressed, the MWF leverages the statistical properties of the sky components, which are encoded in their empirical power spectra. The auto- and cross-frequency power spectra, shown in Figure \ref{fig:image1}, form the core inputs for the filter. The auto-spectra (left panel) again highlight the dominance of the CIB at high frequencies. Crucially, the cross-power spectra (middle panel) reveal the strong spatial correlation of the CIB between channels, particularly between 353, 545, and 857 GHz. It is this statistical correlation that the MWF exploits to distinguish the CIB from other components. The right panel of Figure \ref{fig:image1} explicitly shows that the ground-truth tSZ signal is sub-dominant to the total power in every frequency channel across all angular scales, underscoring the difficulty of the signal extraction task.

\begin{figure}[h!]
    \centering
    \includegraphics[width=\textwidth]{../input_files/plots/step_2_covariance_spectra_2_20260416_235303.png}
    \caption{The empirical auto- and cross-power spectra used as inputs for the Multi-Frequency Wiener Filter. \textit{Left}: The auto-spectra ($C_\ell$) of the six frequency channels, illustrating the extreme dynamic range where the Cosmic Infrared Background (CIB) dominates the high-frequency channels (545 and 857 GHz) by several orders of magnitude. \textit{Middle}: The absolute cross-power spectra between selected frequency pairs, highlighting the strong spatial correlation of the CIB between the 353 and 857 GHz channels, which is leveraged by the filter to model this foreground. \textit{Right}: The auto-power spectrum of the ground-truth tSZ signal (black) and its cross-correlation with individual frequency channels. A comparison with the left panel demonstrates that the target tSZ signal is orders of magnitude fainter than the total power in any given frequency map, underscoring the signal extraction challenge.}
    \label{fig:image1}
\end{figure}

\subsection{Harmonic-space reconstruction fidelity}

We first evaluate the quality of the reconstructed Compton-$y$ maps in the harmonic domain. Figure \ref{fig:image3} compares the performance of the MWF and ILC methods using the transfer function, $T(\ell)$, and the cross-correlation coefficient, $r_\ell$. The transfer function measures the fractional power of the true signal recovered in the reconstruction. As seen in the left panel, the ILC method exhibits a transfer function greater than unity at intermediate scales ($1000 < \ell < 3000$), indicating that foreground residuals and noise are adding excess power. At high multipoles ($\ell > 3000$), the ILC significantly amplifies noise, a known consequence of forcing a foreground null in noisy channels.

In contrast, the MWF demonstrates the behavior of an optimal linear filter. Its transfer function remains close to unity at large scales where the signal-to-noise ratio (SNR) is high and then smoothly rolls off to zero at small scales where noise dominates. This roll-off is not a failure of the method but rather its intended function: by suppressing power in low-SNR modes, the MWF minimizes the mean-squared error of the reconstruction. The cross-correlation coefficient, $r_\ell$ (right panel), shows that both methods maintain excellent phase coherence with the true signal for $\ell < 2000$, after which the correlation drops as noise begins to dominate.

\begin{figure}[h!]
    \centering
    \includegraphics[width=0.9\textwidth]{../input_files/plots/step_4_reconstruction_eval_4.png}
    \caption{Harmonic-space performance of the Multi-Frequency Wiener Filter (MWF) versus an Internal Linear Combination (ILC) baseline, showing the transfer function $T(\ell)$ (left) and the cross-correlation coefficient $r_\ell$ (right). The ILC reconstruction suffers from excess power ($T(\ell) > 1$) and noise amplification at intermediate and high multipoles, respectively. In contrast, the MWF provides a stable reconstruction by optimally suppressing noise at high $\ell$ where the signal-to-noise ratio is low. Both methods achieve a high cross-correlation ($r_\ell \approx 1$) with the true signal at low multipoles ($\ell < 2000$), which then drops to zero as noise begins to dominate the signal.}
    \label{fig:image3}
\end{figure}

\subsection{Galaxy cluster detection and photometry}

To assess the impact of the reconstruction quality on a key scientific application, we performed cluster detection on the reconstructed $y$-maps. This allows for a direct comparison of the methods in terms of catalog completeness, purity, and photometric accuracy.

\subsubsection{Cluster detection pipeline and templates}
Cluster candidates were identified using a spatial matched filter. The templates for this filter were derived from the beam-convolved, radially averaged profiles of simulated clusters, as shown in Figure \ref{fig:image2}. This ensures that the search template matches the expected signal morphology in the final maps, which are smoothed by both the instrumental beams and the filtering process itself. The smoothing effect is most pronounced for the Planck HFI channels, which have larger beams.

\begin{figure}[h!]
    \centering
    \includegraphics[width=0.9\textwidth]{../input_files/plots/step_3_cluster_profiles_3_20260416_235510.png}
    \caption{Spatial templates used for the matched-filter cluster detection pipeline. The right panels show the two-dimensional cluster profile for each frequency channel, convolved with the corresponding instrumental beam, while the left panel displays their radially averaged profiles against the original, unconvolved profile (dashed black line). The smoothing effect and peak suppression are most pronounced for the Planck HFI channels (353, 545, and 857 GHz) due to their larger beams. Using these frequency-specific, beam-convolved templates is essential for the matched filter to accurately identify clusters in the reconstructed maps.}
    \label{fig:image2}
\end{figure}

\subsubsection{Completeness and purity}
The results of the detection pipeline are summarized in Figure \ref{fig:image4}. While both methods achieve similar completeness for the most massive clusters (top left panel), their performance in terms of purity is starkly different. The MWF-derived catalog maintains a high purity ($>80\%$) across a wide range of recovered integrated-$Y$ values. In contrast, the purity of the ILC catalog drops dramatically for all but the most significant detections, indicating that it is dominated by false positives.

The bottom panels of Figure \ref{fig:image4} reveal the source of this discrepancy. The purity of the ILC catalog is strongly anti-correlated with the local intensity of the CIB, degrading significantly in regions with bright CIB emission. This demonstrates that the ILC method fails to adequately separate the CIB, leading to spurious detections where CIB fluctuations are misidentified as clusters. The MWF, by explicitly modeling the CIB as a source of correlated noise, effectively suppresses these fluctuations, resulting in a purity that is largely independent of the CIB environment. This highlights a key advantage of the MWF: robust performance against spatially correlated foregrounds.

\begin{figure}[h!]
    \centering
    \includegraphics[width=0.9\textwidth]{../input_files/plots/step_6_completeness_purity_6_20260417_000631.png}
    \caption{Comparison of cluster detection performance for the Multi-Frequency Wiener Filter (MWF) and an Internal Linear Combination (ILC) baseline. While both methods achieve similar completeness for high-mass clusters (top left), the MWF provides substantially higher purity across the full range of recovered Compton-Y values (top right), whereas the ILC catalog is dominated by false positives at lower signal-to-noise. The bottom panels demonstrate that this performance difference is driven by robustness to Cosmic Infrared Background (CIB) contamination. The MWF maintains high purity even in regions of intense CIB emission, while the ILC's purity degrades significantly, indicating that the ILC misidentifies CIB fluctuations as clusters.}
    \label{fig:image4}
\end{figure}

\subsubsection{Photometric accuracy and bias}
Finally, we examine the quality of the cluster photometry by comparing the recovered integrated Compton-$Y$ for correctly identified clusters with their true values. Figure \ref{fig:image5} shows the resulting mass-observable relation and the fractional bias.

The left panel demonstrates that the MWF produces a significantly tighter relation with lower scatter compared to the ILC. The ILC reconstruction suffers from large upward scatter, particularly for low-mass clusters, which is a direct consequence of positive CIB residuals biasing the recovered flux.

The right panel quantifies the bias in the recovered photometry. The ILC exhibits a significant positive bias, confirming that CIB contamination systematically inflates the measured cluster fluxes. The MWF, on the other hand, shows a consistent negative bias of approximately 20-30\%. This suppression is a characteristic feature of the Wiener filter, which optimally attenuates the signal by a factor of $S/(S+N)$ to minimize the total error. While this introduces a predictable bias that must be accounted for in cosmological analyses (e.g., through calibration with simulations), it results in a measurement with substantially lower variance and higher purity, making it a more robust observable for precision cosmology.

\begin{figure}[h!]
    \centering
    \includegraphics[width=0.9\textwidth]{../input_files/plots/step_6_scatter_bias_6_20260417_000631.png}
    \caption{Comparison of the recovered integrated Compton-$Y$ and fractional mass bias for the Multi-Frequency Wiener Filter (MWF, blue) and the Internal Linear Combination (ILC, red). The left panel shows that the MWF yields a tighter relation between recovered and true integrated $Y$ with lower scatter than the ILC. The right panel quantifies the fractional mass bias, revealing a large positive bias for the ILC due to Cosmic Infrared Background (CIB) contamination. In contrast, the MWF exhibits a consistent negative bias, a characteristic suppression from the Wiener filter that minimizes the mean-square error.}
    \label{fig:image5}
\end{figure}

\end{document}
                